% Part: turing-machines
% Chapter: undecidability
% Section: enumerating-tms

\documentclass[../../../include/open-logic-section]{subfiles}

\begin{document}

\olfileid{tur}{und}{enu}
\olsection{تعداد آلات تورنغ}

\begin{explain}
نبين أن مجموعة آلات تورنغ !!{enumerable}، على معنى ستأتي دقته.
ووجه ذلك أن لكل آلة وصفًا منتهيًا: فحالاتها منتهية، وكذلك أبجدية
شريطها، ودالة انتقالها دالة جزئية من $\OLId{latin-looped}{Q} \times \Sigma$ إلى
$\OLId{latin-looped}{Q} \times \Sigma \times \{\TMleft,\TMright,\TMstay\}$. فيكفي في تعيينها
سرد قيمها عند أزواج الوسيطات التي تُعرّف عليها، وهي منتهية العدد.

إلا أن في هذا القول تفصيلًا ينبغي التنبه له. فإن تعريف آلة تورنغ
لم يقيّد ال!!{element}s التي تُتخذ حالات أو رموزًا للشريط؛ فيجوز،
بحسب ظاهر التعريف، أن نجعل أي عدد حقيقي بعينه حالة لآلة ما. وإنما
العبرة ب\emph{سلوك} الآلة، لا بذوات الأشياء المسماة حالات ورموزًا.
ولننظر في الآلتين المرسومتين في \olref{fig:variants}.
\begin{figure}
\begin{center}
\begin{tikzpicture}[->,>=stealth',shorten >=1pt,auto,node distance=2.8cm,
                    semithick]
  \tikzstyle{every state}=[fill=none,draw=black,text=black]

  \node[initial,state]         (A)              {$\OLId{latin-ordinary}{q}_\OLMathNumeral{0}$};
  \node[state]         (B) [right of=A] {$\OLId{latin-ordinary}{q}_\OLMathNumeral{1}$};

  \path (A) edge [bend left] node {\TMtrans{\TMstroke}{\TMstroke}{\TMright}} (B)
        (B) edge [loop above] node {\TMtrans{\TMblank}{\TMblank}{\TMright}} (B)
            edge [bend left] node {\TMtrans{\TMstroke}{\TMstroke}{\TMright}} (A);
\end{tikzpicture}\\
\begin{tikzpicture}[->,>=stealth',shorten >=1pt,auto,node distance=2.8cm,
  semithick]
\tikzstyle{every state}=[fill=none,draw=black,text=black]

\node[initial,state]         (A)              {$\OLId{latin-ordinary}{s}$};
\node[state]         (B) [right of=A] {$\OLId{latin-ordinary}{h}$};

\path (A) edge [bend left] node {\TMtrans{A}{A}{\TMright}} (B)
(B) edge [loop above] node {\TMtrans{\TMblank}{\TMblank}{\TMright}} (B)
edge [bend left] node {\TMtrans{A}{A}{\TMright}} (A);
\end{tikzpicture}
\end{center}
\caption{اختلاف التسميات في آلة \emph{الزوجية}}
\ollabel{fig:variants}
\end{figure}
فالمخطط الأول لآلة $\OLId{latin-looped}{M}$ أبجدية شريطها
$\Sigma=\{\TMendtape,\TMblank,\TMstroke\}$ وحالاتها
$\{\OLId{latin-ordinary}{q}_\OLMathNumeral{0},\OLId{latin-ordinary}{q}_\OLMathNumeral{1}\}$، والثاني لآلة $\OLId{latin-looped}{M}'$ أبجديتها
$\Sigma'=\{\TMendtape,\TMblank,\OLId{latin-looped}{A}\}$ وحالتاها $\{\OLId{latin-ordinary}{s},\OLId{latin-ordinary}{h}\}$. ولا فرق بين
تعليماتهما سوى هذه التسميات: فالآلة $\OLId{latin-looped}{M}$ تتوقف عند متتالية من $\OLId{latin-ordinary}{n}$
علامات~$\TMstroke$ إذا وفقط إذا كان $\OLId{latin-ordinary}{n}$ زوجيًا؛ وكذلك تتوقف $\OLId{latin-looped}{M}'$
عند متتالية من $\OLId{latin-ordinary}{n}$ رموز~$\OLId{latin-looped}{A}$ إذا وفقط إذا كان $\OLId{latin-ordinary}{n}$ زوجيًا. فليس
التغيير إلا تسمية $\TMstroke$ باسم~$\OLId{latin-looped}{A}$، و$\OLId{latin-ordinary}{q}_\OLMathNumeral{0}$ باسم~$\OLId{latin-ordinary}{s}$، و$\OLId{latin-ordinary}{q}_\OLMathNumeral{1}$
باسم~$\OLId{latin-ordinary}{h}$. ويجري هذا الحكم على سائر الآلات: نعد الآلتين متماثلتين
متى أمكن الانتقال من إحداهما إلى الأخرى بإعادة تسمية الحالات والرموز.
بل يجوز أن تكون أسماؤها كلها أعدادًا صحيحة موجبة؛ فنضع~$\OLMathNumeral{1}$
موضع~$\sigma_\OLMathNumeral{0}$، و~$\OLMathNumeral{2}$ موضع~$\sigma_\OLMathNumeral{1}$، وهكذا. وتكون تسمية
$\TMendtape$ هي~$\OLMathNumeral{1}$، وتسمية $\TMblank$ هي~$\OLMathNumeral{2}$، وعلى هذا القياس.
فتصير آلة \emph{الزوجية} إلى الصورة التي في
\olref{fig:standard-even}.
\begin{figure}
\[\begin{tikzpicture}[->,>=stealth',shorten >=\OLMathNumeral{1}pt,auto,node distance=\OLMathNumeral{2}.\OLMathNumeral{8}cm,
  semithick]
\tikzstyle{every state}=[fill=none,draw=black,text=black]

\node[initial,state]         (\OLId{latin-looped}{A})              {$\OLMathNumeral{1}$};
\node[state]         (\OLId{latin-looped}{B}) [right of=\OLId{latin-looped}{A}] {$\OLMathNumeral{2}$};

\path (\OLId{latin-looped}{A}) edge [bend left] node {\TMtrans{\OLMathNumeral{3}}{\OLMathNumeral{3}}{\TMright}} (\OLId{latin-looped}{B})
(\OLId{latin-looped}{B}) edge [loop above] node {\TMtrans{\OLMathNumeral{2}}{\OLMathNumeral{2}}{\TMright}} (\OLId{latin-looped}{B})
edge [bend left] node {\TMtrans{\OLMathNumeral{3}}{\OLMathNumeral{3}}{\TMright}} (\OLId{latin-looped}{A});
\end{tikzpicture}
\]
\caption{آلة \emph{زوجية} قياسية}
\ollabel{fig:standard-even}
\end{figure}
ولنسمّ آلة تورنغ \emph{قياسية} إذا كانت حالاتُها ورموزُها أعدادًا
صحيحة موجبة، ولنقصر النظر فيما يأتي على هذه الآلات.
\footnote{«الآلة القياسية» اصطلاح نستعمله هنا، وليس اصطلاحًا متعارفًا عليه.}

كان المطلوب أن مجموعة آلات تورنغ !!{enumerable}؛ وقد تبين أن
حسبنا لهذا الغرض إثبات أن مجموعة الآلات القياسية !!{enumerable}.
ولتكن $\OLId{latin-looped}{M}=\tuple{\OLId{latin-looped}{Q},\Sigma,\OLId{latin-ordinary}{q}_\OLMathNumeral{0},\delta}$ آلة قياسية. وصفها بسلسلة
منتهية من الأعداد الصحيحة الموجبة يكون على الوجه الآتي: نسرد عدد
الحالات، فالحالات، ثم عدد الرموز، فالرموز، ثم الحالة الابتدائية.
وكل ذلك أعداد صحيحة موجبة، لأن $\OLId{latin-looped}{M}$ قياسية. ويبقى تعيين~$\delta$.
وأزواج الوسيطات $\tuple{\OLId{latin-ordinary}{q},\sigma}$ منتهية العدد، لانتهاء $\OLId{latin-looped}{Q}$
و~$\Sigma$؛ فلا تحمل~$\delta$ إلا قائمة منتهية من الخماسيات
$\tuple{\OLId{latin-ordinary}{q},\sigma,\OLId{latin-ordinary}{q}',\sigma',\OLId{latin-ordinary}{d}}$ التي يتحقق فيها
$\delta(\OLId{latin-ordinary}{q},\sigma)=\tuple{\OLId{latin-ordinary}{q}',\sigma',\OLId{latin-looped}{D}}$، ويكون $\OLId{latin-ordinary}{d}$ رمزًا عدديًا
للاتجاه~$\OLId{latin-looped}{D}$: نجعل $\OLMathNumeral{1}$ رمز~$\TMleft$، و$\OLMathNumeral{2}$ رمز~$\TMright$،
و$\OLMathNumeral{3}$ رمز~$\TMstay$.

فقد صار وصف كل آلة قياسية قائمة منتهية من الأعداد الصحيحة الموجبة،
أي متتالية $\OLId{latin-ordinary}{s}_\OLId{latin-looped}{M}\in(\PosInt)^*$. ووصف آلة \emph{الزوجية} القياسية،
مثلًا، هو
\[
\OLMathNumeral{2}, \underbrace{\OLMathNumeral{1}, \OLMathNumeral{2}}_\OLId{latin-looped}{Q}, \OLMathNumeral{3}, \overbrace{\OLMathNumeral{1}, \OLMathNumeral{2}, \OLMathNumeral{3}}^\Sigma, \OLMathNumeral{1}, \underbrace{\OLMathNumeral{1}, \OLMathNumeral{3}, \OLMathNumeral{2}, \OLMathNumeral{3}, \OLMathNumeral{2}}_{\delta(\OLMathNumeral{1},\OLMathNumeral{3}) = \tuple{\OLMathNumeral{2},\OLMathNumeral{3},\OLId{latin-looped}{R}}},
\overbrace{\OLMathNumeral{2}, \OLMathNumeral{2}, \OLMathNumeral{2}, \OLMathNumeral{2}, \OLMathNumeral{2}}^{\delta(\OLMathNumeral{2},\OLMathNumeral{2}) = \tuple{\OLMathNumeral{2},\OLMathNumeral{2},\OLId{latin-looped}{R}}},
\underbrace{\OLMathNumeral{2}, \OLMathNumeral{3}, \OLMathNumeral{1}, \OLMathNumeral{3}, \OLMathNumeral{2}}_{\delta(\OLMathNumeral{2},\OLMathNumeral{3}) = \tuple{\OLMathNumeral{1},\OLMathNumeral{3},\OLId{latin-looped}{R}}}.
\]
\end{explain}

\begin{thm}
ليست كل دالة من $\Nat$ إلى~$\Nat$ قابلة للحساب بآلة تورنغ.
\end{thm}

\begin{proof}
مجموعة المتتاليات المنتهية من الأعداد الصحيحة الموجبة
$(\PosInt)^*$ !!{enumerable}
(\cref{sfr:siz:zigzag:prob:posint-star}). وأوصاف الآلات القياسية
مجموعة جزئية من~$(\PosInt)^*$، فتقبل التعداد أيضًا. وكل دالة من
$\Nat$ إلى~$\Nat$ قابلة للحساب بآلة تورنغ تحسبها آلة، بل آلات كثيرة.
فإذا سمينا حالات إحدى هذه الآلات ورموزها بأعداد صحيحة موجبة، وجعلنا
خاصةً $\TMendtape$ هو~$\OLMathNumeral{1}$، و$\TMblank$ هو~$\OLMathNumeral{2}$،
و$\TMstroke$ هو~$\OLMathNumeral{3}$، حصلنا على آلة قياسية تحسب الدالة نفسها.
إذن كل دالة قابلة للحساب تحسبها آلة قياسية، ومن ثم تقبل مجموعة
الدوال القابلة للحساب من $\Nat$ إلى~$\Nat$ التعداد.

لكن مجموعة الدوال كلها من $\Nat$ إلى~$\Nat$ ليست !!{enumerable}
(\cref{sfr:siz:red:prob:nat-nat}). فلو كانت كل دالة تحسبها آلة
تورنغ، لكان سرد أوصاف الآلات التي تحسبها تعدادًا للدوال جميعًا، وهو
خلاف ما ثبت. فلا بد من دوال لا تقبل الحساب بآلة تورنغ.
\end{proof}

\begin{prob}
  التمس وصفًا لآلات تورنغ يستغني عن سرد الحالات ورموز الأبجدية واحدًا
  واحدًا. ولك أن تختار تعريفًا آخر للآلة «القياسية»؛ على أن تبين
  كيف تستطيع آلة قياسية، بهذا المعنى الذي اخترته، محاكاة كل آلة تورنغ.
\end{prob}

\end{document}
